\newcommand{\tipoRequerimento}{Cobrança indevida sobre enquadramento tarifário incorreto}
\section{DOS FATOS E DOS DIREITOS}

% Contestação geral para cargas presentes que divergem da exclusividade de circuitos de iluminação

\begin{enumerate}[resume]

\item Primeiramente, ainda que a unidade consumidora possua tomadas, pontos de conexão ou extensões sem equipamentos conectados no local, tal circunstância, por si só, não descaracteriza a natureza predominantemente pública da instalação nem afasta seu enquadramento na subclasse de Iluminação Pública (IP).

\item Com efeito, a existência de tomadas ou extensões deve ser analisada em conjunto com a finalidade efetivamente atribuída à unidade consumidora, não sendo suficiente a mera constatação de cargas ou pontos de utilização acessórios para justificar, de forma automática, sua classificação como unidade pertencente ao Poder Público.

\item Nesse sentido, quando a instalação elétrica estiver destinada primordialmente à iluminação de quadras esportivas, campos de futebol ou outros espaços públicos de lazer de livre acesso, a eventual presença de tomadas ou extensões voltadas ao apoio, manutenção ou utilização desses espaços não altera a finalidade essencial do fornecimento, que permanece vinculada à iluminação de área pública.

\item Contudo, a mera existência de tomadas, extensões ou outros pontos de fornecimento, desacompanhada da demonstração de que tais estruturas constituem utilização autônoma e predominante de energia para finalidade diversa da iluminação pública, não fornece fundamento suficiente para afastar o enquadramento na subclasse de Iluminação Pública (IP). Do mesmo modo, a simples invocação genérica do artigo 189 da Resolução Normativa nº 1.000/2021 não supre a necessidade de demonstração concreta dos pressupostos que justificariam a classificação adotada pela concessionária.

\item Desse modo, se as tomadas ou extensões existentes possuem caráter meramente acessório e estão vinculadas à utilização ou manutenção do espaço público iluminado, não se mostra juridicamente adequada a alteração da natureza da unidade consumidora para a classificação de Poder Público. Para tanto, seria indispensável comprovar que a unidade efetivamente possui utilização independente e relevante para finalidade diversa daquela própria da iluminação pública.

\item Por fim, inexistindo demonstração concreta de que a presença de tomadas ou extensões tenha descaracterizado a finalidade de iluminação pública, permanece indevida a cobrança decorrente do enquadramento incorreto da unidade consumidora, assegurando-se ao consumidor o direito à restituição dos valores cobrados e pagos indevidamente, na forma prevista pela regulamentação aplicável.

\end{enumerate}


\section{DOS PEDIDOS}

\begin{enumerate}[resume]
\item 
Diante do exposto, requer que a Concessionária:

\begin{enumerate}
\item 
Proceda com a correção do enquadramento tarifário da mesma em conformidade com o estabelecido na Resolução Normativa Nº 1.000/2021 da ANEEL.
\item 
Efetue a devolução de todos os valores cobrados a maior, referente aos últimos 10 (dez) anos. Conforme o Despacho Nº 2.006 da ANEEL de 08 de julho de 2024, observando o prazo prescricional de 10 anos, conforme previsto no art. 205 do Código Civil.
\item 
Devolva os valores cobrados a maior em dobro, atualizados monetariamente pelo IPCA e juros de mora de 1\% ao mês calculados pro \textit{rata die}, conformidade com o estabelecido na Resolução Normativa Nº 1.000/2021 da ANEEL.   
\item 
Que apresente as seguintes informações: Data de ligação da UC, se houve modificação de tarifa entre a data da ligação e a data da resposta.   
\item 
Que em caso de indeferimento apresente os argumentos e as provas (comprovações) para não enquadramento (fotos, carga instalada etc).   
\item 
Proceda com a devolução do valor integral, por meio de depósito bancário, na conta corrente do Município, conforme estabelece o § 7º, do art. 323 da REN Nº 1.000/2021 da ANEEL, comunicando por escrito o resultado do deferimento desta reclamação, por meio de e-mail.

\end{enumerate}
\end{enumerate}